%!TEX root = ./main.tex

\section[Hybrid]{Version 3: Meaning-Aware Bits}

% SECTION TIMING: 39:40--43:10 (2 frames): 100 + 110 = 210 s.

% Teaching note (100 s): Version 3 keeps today's pipeline -- standard codec, standard
% network -- and changes only WHERE the bits go. An AI on the camera marks what
% matters for the task; the standard encoder (JPEG/H.264) spends more bits there and
% fewer on the rest; an AI at the receiver fills in what was dropped or lost. It does
% not climb Weaver's levels: it keeps Level A's pipe and uses Level B/C knowledge to
% decide which bits matter most. Tie back to "Same loss, two outcomes" (the
% pedestrian deserves more protection than the sky) and to "Where it struggles"
% ("every bit counts the same"). Why it matters for an edge-networks course: it needs
% nothing new in radios, codecs, or protocols, so it is the easiest kind to deploy.
% Networking cousins (verified earlier): AccMPEG (MLSys 2022) gives higher quality to
% the macroblocks that matter for accuracy; DDS (ACM SIGCOMM 2020) sends a
% low-quality stream and re-sends important regions in high quality; MPEG's Video
% Coding for Machines (ISO/IEC 23888-2) codes video for machine analysis.
\begin{frame}{Version 3: keep the pipe, spend bits by meaning}
  \vspace{1.0em}
  \centering
  \begin{tikzpicture}[scale=0.68,transform shape,
      std/.style={netnode,minimum width=2.2cm,minimum height=1.0cm,font=\small\bfseries},
      ai/.style={nnnode,minimum width=2.2cm,minimum height=1.0cm,font=\small\bfseries}]
    \node[endnode] (v) at (0,0) {Video};
    \node[ai] (score) at (2.75,0) {AI: what\\matters?};
    \node[std] (enc) at (5.6,0) {H.264\\[-2pt]{\scriptsize\mdseries standard}};
    \node[chnode,minimum width=2.2cm,minimum height=1.0cm,font=\small\bfseries] (net)
      at (8.45,0) {Today's\\network};
    \node[ai] (fill) at (11.3,0) {AI fills\\in gaps};
    \node[endnode] (rx) at (13.85,0) {Receiver};
    \draw[flow] (v) -- (score); \draw[flow] (score) -- (enc); \draw[flow] (enc) -- (net);
    \draw[flow] (net) -- (fill); \draw[flow] (fill) -- (rx);
    \node[lbl,below=3pt of enc] {unchanged};
    \node[lbl,below=3pt of net] {unchanged};
  \end{tikzpicture}

  \vspace{0.9em}
  \begin{tikzpicture}[scale=0.85,visible on=<2->]
    \pic at (0,0) {streetv3};
    \draw[-{Latex[length=2mm]},softgreen,thick] (4.3,1.55) -- (1.55,1.3);
    \node[anchor=west,font=\small\bfseries,text=softgreen] at (4.35,1.55)
      {many bits: the pedestrian};
    \draw[-{Latex[length=2mm]},softgray,thick] (4.3,0.55) -- (2.9,1.6);
    \node[anchor=west,font=\small\bfseries,text=softgray] at (4.35,0.55)
      {few bits: sky, road, car};
  \end{tikzpicture}

  \vspace{0.7em}
  \uncover<3->{{\large Nothing new for radios, codecs, or networks.}\par}

  \vspace{0.7em}
  \takeawayat{4-}{Same pipe, smarter spending ---\\
    the easiest version to deploy today.}
\end{frame}

% Teaching note (110 s): one real system, so students see Version 3 works. KDC
% (X. Chen, Z. Feng, W. Zhao, J. Ding, K. Sun, X. Zhang, "From Bits to Tokens:
% Knowledge-Driven Generative Communication of Multimodal Data," USENIX NSDI 2026;
% read in full): a foundation model (CLIP, plus a small language model run once
% offline) turns the task description into keywords -- "autonomous driving" -> car,
% pedestrian, traffic light -- and scores every image patch against them; the
% standard JPEG/H.264 encoder then keeps finer detail in high-scoring regions and
% quantizes the rest coarsely (Sec. 3.2: "fully compatible with H.264/JPEG/PCC and
% existing hardware accelerators"; it runs at the application layer, PHY/MAC and
% retransmissions unchanged); the keywords travel once in a small header; a
% receiver-side restoration network (UNet conditioned on a knowledge base) fills in
% lost or coarse parts. Numbers (Sec. 5.1, images: ImageNet-1K and Cityscapes, 5G
% traces): ~0.85 Top-1 with only 10 kB, close to ~0.87 for the full image, where a
% full JPEG averages ~30 kB; 0.55-0.65 Top-1 at every packet-loss level, while
% progressive JPEG and the learned codec LimitNet fall to <=0.2 at 5-10% loss and the
% semantic system StarFish to ~0.3 at 10% loss (plain JPEG fails). Video (Sec. 5.2,
% driving-video segmentation over real 5G traces): KDC on top of H.264 keeps the
% segmentation score (IoU) 0.1-0.2 higher than plain H.264 under medium-to-high loss.
% Weak spots: it still sends a picture (more bits than Version 2), and it needs AI
% on both ends (the paper reports 32 frames/s on a Jetson Orin).
\begin{frame}{Version 3 in practice}
  \vspace{0.9em}
  \centering
  {\large Example: \textbf{KDC} (NSDI 2026)\par}

  \vspace{0.8em}
  \begin{tikzpicture}[scale=0.71,transform shape,
      st/.style={nnnode,minimum width=3.4cm,minimum height=1.2cm,font=\small\bfseries}]
    \node[st] (a) at (0,0) {Task $\to$ keywords\\[-2pt]{\scriptsize\mdseries ``car, pedestrian''}};
    \node[st,draw=ranblue,fill=blue!5] (b) at (4.3,0)
      {JPEG / H.264\\[-2pt]{\scriptsize\mdseries more detail on those}};
    \node[st] (c) at (8.6,0) {Receiver AI\\[-2pt]{\scriptsize\mdseries fills in what was lost}};
    \draw[flow] (a) -- (b); \draw[flow] (b) -- (c);
  \end{tikzpicture}

  \vspace{0.9em}
  \raggedright
  \begin{itemize}[<+(1)->]\itemsep5pt
    \item Images: near-full accuracy with \textbf{10 kB} \textcolor{softgray}{(a full
      JPEG $\approx$30 kB)}
    \item 10\% packet loss: \textbf{0.55--0.65} accuracy vs.\ \textbf{$\le$0.3} for others
    \item Driving video: H.264 \textbf{+ KDC} beats plain H.264 under loss
  \end{itemize}

  \vspace{0.6em}
  \takeawayat{5-}{Works with today's codecs and networks ---\\
    but still sends a picture, and needs AI on both ends.}
  \source{https://www.usenix.org/conference/nsdi26/presentation/chen-xingyu}{X.\ Chen
    et al., ``From Bits to Tokens: Knowledge-Driven Generative Communication,''
    USENIX NSDI, 2026}
\end{frame}
